% Options for packages loaded elsewhere
\PassOptionsToPackage{unicode}{hyperref}
\PassOptionsToPackage{hyphens}{url}
\documentclass[
  american,
  11pt,
  a4paper,
]{article}
\usepackage{xcolor}
\usepackage[margin=26mm]{geometry}
\usepackage{amsmath,amssymb}
\setcounter{secnumdepth}{5}
\usepackage{iftex}
\ifPDFTeX
  \usepackage[T1]{fontenc}
  \usepackage[utf8]{inputenc}
  \usepackage{textcomp} % provide euro and other symbols
\else % if luatex or xetex
  \usepackage{unicode-math} % this also loads fontspec
  \defaultfontfeatures{Scale=MatchLowercase}
  \defaultfontfeatures[\rmfamily]{Ligatures=TeX,Scale=1}
\fi
\usepackage{lmodern}
\ifPDFTeX\else
  % xetex/luatex font selection
  \setmainfont[BoldFont=lmroman10-bold.otf,ItalicFont=lmroman10-italic.otf,BoldItalicFont=lmroman10-bolditalic.otf]{lmroman10-regular.otf}
  \setsansfont[BoldFont=lmsans10-bold.otf,ItalicFont=lmsans10-oblique.otf,BoldItalicFont=lmsans10-boldoblique.otf]{lmsans10-regular.otf}
  \setmonofont[ItalicFont=lmmono10-italic.otf]{lmmono10-regular.otf}
  \setmathfont[]{latinmodern-math.otf}
\fi
% Use upquote if available, for straight quotes in verbatim environments
\IfFileExists{upquote.sty}{\usepackage{upquote}}{}
\IfFileExists{microtype.sty}{% use microtype if available
  \usepackage[]{microtype}
  \UseMicrotypeSet[protrusion]{basicmath} % disable protrusion for tt fonts
}{}
\usepackage{setspace}
\makeatletter
\@ifundefined{KOMAClassName}{% if non-KOMA class
  \IfFileExists{parskip.sty}{%
    \usepackage{parskip}
  }{% else
    \setlength{\parindent}{0pt}
    \setlength{\parskip}{6pt plus 2pt minus 1pt}}
}{% if KOMA class
  \KOMAoptions{parskip=half}}
\makeatother
\usepackage{longtable,booktabs,array}
\usepackage{calc} % for calculating minipage widths
% Correct order of tables after \paragraph or \subparagraph
\usepackage{etoolbox}
\makeatletter
\patchcmd\longtable{\par}{\if@noskipsec\mbox{}\fi\par}{}{}
\makeatother
% Allow footnotes in longtable head/foot
\IfFileExists{footnotehyper.sty}{\usepackage{footnotehyper}}{\usepackage{footnote}}
\makesavenoteenv{longtable}
\ifLuaTeX
\usepackage[bidi=basic,shorthands=off,]{babel}
\else
\usepackage[bidi=default,shorthands=off,]{babel}
\fi
\ifPDFTeX
\else
\babelfont{rm}[BoldFont=lmroman10-bold.otf,ItalicFont=lmroman10-italic.otf,BoldItalicFont=lmroman10-bolditalic.otf]{lmroman10-regular.otf}
\fi
\ifLuaTeX
  \usepackage{selnolig} % disable illegal ligatures
\fi
\setlength{\emergencystretch}{3em} % prevent overfull lines
\providecommand{\tightlist}{%
  \setlength{\itemsep}{0pt}\setlength{\parskip}{0pt}}
% Typography for the reproducible mathematics guide.
\usepackage{titlesec}
\usepackage{fancyhdr}
\usepackage{needspace}
\usepackage{etoolbox}
\definecolor{GuideInk}{HTML}{21343C}
\definecolor{GuideAccent}{HTML}{176B61}
\definecolor{GuideMuted}{HTML}{65747C}
\definecolor{GuidePale}{HTML}{EDF3F1}

\setlength{\headheight}{15pt}
\setlength{\headsep}{18pt}
\setlength{\footskip}{28pt}
\pagestyle{fancy}
\fancyhf{}
\fancyhead[L]{\sffamily\footnotesize\textcolor{GuideMuted}{HASSETT LATTICE EXPLORER}}
\fancyhead[R]{\sffamily\footnotesize\textcolor{GuideMuted}{MATHEMATICS GUIDE}}
\fancyfoot[L]{\sffamily\scriptsize\textcolor{GuideMuted}{Exact arithmetic\quad /\quad November 2026}}
\fancyfoot[R]{\sffamily\small\textcolor{GuideInk}{\thepage}}
\renewcommand{\headrulewidth}{0.35pt}
\renewcommand{\footrulewidth}{0pt}

\titleformat{\section}{\sffamily\Large\bfseries\color{GuideInk}}{\thesection}{0.7em}{}
\titleformat{\subsection}{\sffamily\large\bfseries\color{GuideAccent}}{\thesubsection}{0.7em}{}
\titleformat{\subsubsection}{\sffamily\normalsize\bfseries\color{GuideInk}}{\thesubsubsection}{0.7em}{}
\titlespacing*{\section}{0pt}{21pt plus 3pt minus 2pt}{9pt}
\titlespacing*{\subsection}{0pt}{15pt plus 2pt minus 2pt}{6pt}
\pretocmd{\subsection}{\Needspace{8\baselineskip}}{}{}
\setlength{\parindent}{0pt}
\setlength{\parskip}{5pt plus 1pt minus 1pt}
\setlength{\emergencystretch}{2em}
\setlength{\tabcolsep}{5pt}
\renewcommand{\arraystretch}{1.18}
\clubpenalty=10000
\widowpenalty=10000
\displaywidowpenalty=10000
\predisplaypenalty=10000
\raggedbottom
\AtBeginDocument{\hypersetup{colorlinks=true,linkcolor=GuideAccent,urlcolor=GuideAccent,citecolor=GuideAccent}}

\providecommand{\GuideSubtitle}{}
\providecommand{\subtitle}[1]{\gdef\GuideSubtitle{#1}}
\makeatletter
\renewcommand{\maketitle}{%
  \begin{titlepage}
  \thispagestyle{empty}
  \vspace*{10mm}
  {\sffamily\bfseries\small\color{GuideAccent}RESEARCH SOFTWARE\quad /\quad MATHEMATICS GUIDE\par}
  \vspace{10mm}
  {\sffamily\bfseries\fontsize{32}{37}\selectfont\color{GuideInk}\@title\par}
  \vspace{8mm}
  {\sffamily\large\color{GuideMuted}\GuideSubtitle\par}
  \vspace{16mm}
  {\color{GuideAccent}\hrule height 1.1pt}
  \vspace{12mm}
  {\large
  \[
  G_t=\begin{pmatrix}3&\epsilon_1&\epsilon_2\\
  \epsilon_1&a_1&t\\\epsilon_2&t&a_2\end{pmatrix},
  \qquad
  \det G_t=\frac{d_1d_2-(3t-\epsilon_1\epsilon_2)^2}{3}.
  \]
  }
  \vspace{10mm}
  \begin{minipage}{0.92\textwidth}
  {\large\color{GuideInk}Two divisor labels determine a finite family of rank-three lattices. Saturation and integral isometry turn that family into the irreducible components.\par}
  \vspace{7mm}
  {\sffamily\small\color{GuideMuted}Exact computation\quad\textbullet\quad Published validation\quad\textbullet\quad K3 associations\par}
  \end{minipage}
  \vfill
  {\sffamily\bfseries\color{GuideInk}\@author\par}
  \vspace{2mm}
  {\sffamily\small\color{GuideMuted}\@date\par}
  \end{titlepage}
  \pagenumbering{roman}
}
\makeatother
\apptocmd{\tableofcontents}{\clearpage\pagenumbering{arabic}}{}{}
\usepackage{bookmark}
\IfFileExists{xurl.sty}{\usepackage{xurl}}{} % add URL line breaks if available
\urlstyle{same}
\hypersetup{
  pdftitle={Hassett divisors: from lattices to components},
  pdfauthor={Hassett Lattice Explorer},
  pdflang={en-US},
  hidelinks,
  pdfcreator={LaTeX via pandoc}}

\title{Hassett divisors: from lattices to components}
\usepackage{etoolbox}
\makeatletter
\providecommand{\subtitle}[1]{% add subtitle to \maketitle
  \apptocmd{\@title}{\par {\large #1 \par}}{}{}
}
\makeatother
\subtitle{An exact algorithm, reproducible computations, and the K3
connection}
\author{Hassett Lattice Explorer}
\date{14 November 2026}

\begin{document}
\maketitle

{
\setcounter{tocdepth}{2}
\tableofcontents
}
\setstretch{1.04}
\section{What this project computes}\label{what-this-project-computes}

Given two distinct nonempty Hassett divisors, this project computes the
irreducible components of their reduced set-theoretic intersection in
the moduli space of smooth complex cubic fourfolds. A component is
represented by a positive definite integral lattice of rank three with a
distinguished element of square three. The program returns Gram
matrices, discriminants, saturation certificates, and the exact K3
admissibility status of the generic algebraic lattice. It does not
compute scheme-theoretic intersection multiplicities, equations for
cubic fourfolds, or a birational parametrization.

The implementation is an arithmetic realization of Yang--Yu's Algorithm
7.6 {[}YY{]}. Its two essential features are exhaustive enumeration of
allowable overlattices and identification up to integral isometry
preserving the distinguished class. Counting positive definite Gram
matrices alone does not solve the problem.

All classification decisions use exact arithmetic. Python supplies a
standalone reference engine; Julia reuses Hecke/Nemo, and Macaulay2
reuses its exact arithmetic and lattice facilities. The published
answers are used in regression tests rather than as a database of
component answers. A computation that exceeds its configured resource
budget raises an error rather than reporting a partial list as complete.

The source release includes a \href{math-audit.md}{mathematical audit}
tracing every geometric hypothesis and an independent enumeration using
general Hermite inclusion matrices. The
\href{benchmark-components.md}{nineteen benchmark dossiers} supply
reproducible basis changes, K3 witnesses, and theorem-specific
rationality statements.

\subsection{The geometric input}\label{the-geometric-input}

Let \(X\subset\mathbb P^5_{\mathbb C}\) be a smooth cubic fourfold.
Write
\[A(X)=H^4(X,\mathbb Z)\cap H^{2,2}(X),\qquad o=h^2.\]

The intersection form on \(A(X)\) is positive definite, and \(o^2=3\).
The element \(o\) is distinguished: its orthogonal complement in
\(A(X)\) is even. In the bases used here, this condition is equivalent
to
\[v^2\equiv(o,v)\pmod2\quad\text{for every }v\in A(X).\]

Indeed, if \(o^\perp\) is even, then \(3v-(o,v)o\in o^\perp\) has square
congruent to \(v^2-(o,v)\) modulo two. The converse is immediate by
taking \((o,v)=0\). Consequently, the condition can be tested on a basis
using the diagonal of a Gram matrix.

A discriminant-\(d\) label is a primitive rank-two sublattice
\(K_d\subset A(X)\) containing \(o\), with determinant \(d\). Here
primitive means that the quotient by \(K_d\) has no torsion. This is a
condition on the integral lattice, not just on its rational span.

Hassett's nonemptiness theorem gives precisely the labels
\[d>6,\qquad d\equiv0,2\pmod6.\]

For such a label, the corresponding divisor \(\mathcal C_d\) is
irreducible {[}H, Theorem 1.0.1{]}. Its dimension is 19 inside the
20-dimensional moduli space. At a very general point of a component of
the intersection of two distinct divisors, \(A(X)\) has rank three.
Higher-rank algebraic lattices occur on further special subloci.

\subsection{Why two labels give a rank-three
problem}\label{why-two-labels-give-a-rank-three-problem}

The two labels share \(o\). Their rational spans cannot coincide: if two
primitive rank-two sublattices of the same ambient lattice have the same
rational span, both equal the intersection of that span with the ambient
lattice, hence are equal and have the same discriminant. Distinct
discriminants therefore give a rank-three sum.

One must then saturate that sum in \(A(X)\). The saturation can be
strictly larger, without changing rank. This is why component
discriminants need not be the determinants of the first matrices one
writes down.

The geometric theorem needed to turn this finite lattice search into a
component computation is Yang--Yu's nonemptiness and irreducibility
theorem {[}YY, Theorem 1.1 / 5.1{]}. A positive definite lattice with a
distinguished square-three element, no vector of square two, and rank
plus discriminant-group length at most 20 defines a nonempty irreducible
locus. For rank three, the last inequality is automatic, since the
discriminant group needs at most three generators. The locus has
codimension two. Their Proposition 7.3 and Corollary 7.5 identify the
intersection components in this way.

\section{Normalized Gram matrices}\label{normalized-gram-matrices}

For each label \(d_i\), put
\[\epsilon_i=\begin{cases}0&d_i\equiv0\pmod6,\\1&d_i\equiv2\pmod6,\end{cases}
\qquad a_i=\frac{d_i+\epsilon_i^2}{3}.\]

The two rank-two lattices can be written in bases \((o,u_i)\) with Gram
matrices
\[\begin{pmatrix}3&\epsilon_i\\\epsilon_i&a_i\end{pmatrix}.\]

Thus their sum has a basis \((o,u_1,u_2)\) with matrix
\[G_t=\begin{pmatrix}
3&\epsilon_1&\epsilon_2\\
\epsilon_1&a_1&t\\
\epsilon_2&t&a_2
\end{pmatrix},\qquad t\in\mathbb Z.\]

The unknown pairing \(t=(u_1,u_2)\) records how the two labels sit
together. Expanding the determinant gives a useful exact search bound:
\[\Delta_t=\det G_t
=\frac{d_1d_2-(3t-\epsilon_1\epsilon_2)^2}{3}.\]

The leading minors are \(3\) and \(d_1\), both positive. Sylvester's
criterion therefore reduces positivity to \(\Delta_t>0\), or
\[|3t-\epsilon_1\epsilon_2|<\sqrt{d_1d_2}.\]

The implementation derives the integer bounds with integer square roots
and strict integer inequalities. If one of the \(\epsilon_i\) vanishes,
changing the sign of its generator changes \(t\) to \(-t\), so only one
sign is necessary. This is a basis symmetry, not a heuristic truncation.

The code orders equal-residue discriminants increasingly. For a mixed
pair, it puts the discriminant congruent to two modulo six first, so
that \((\epsilon_1,\epsilon_2)=(1,0)\). Its reported \(t\) and matrices
use this normalized order; the original request is retained separately
as \texttt{input\_pair} in the JSON. In either case with a vanishing
\(\epsilon_i\), the search takes \(t\ge0\).

\subsection{Excluding forbidden vectors
exactly}\label{excluding-forbidden-vectors-exactly}

For each positive definite matrix, the program searches for a vector of
square two. This search is finite. If \(G\) is positive definite and
\(v^TGv=m\), the Cauchy--Schwarz inequality in the \(G\)-metric gives
\[v_i^2\le m(G^{-1})_{ii}.\]

Since \(G^{-1}=\operatorname{adj}(G)/\det G\), every coordinate bound
can be evaluated using integer arithmetic. The engine uses a sharper
two-variable version of this argument. Write
\(G=\left(\begin{smallmatrix}3&\mathbf p^T\\\mathbf p&B\end{smallmatrix}\right)\)
and \(H=3B-\mathbf p\mathbf p^T\). For \(v=(x,y,z)^T\) with \(v^TGv=m\)
and \((o,v)=k\), one has
\[
\begin{pmatrix}y&z\end{pmatrix}H\begin{pmatrix}y\\z\end{pmatrix}=3m-k^2,
\qquad x=\frac{k-p_1y-p_2z}{3}.
\]

For square-two vectors, \(k^2\le6\), so only \(k=-2,-1,0,1,2\) occur.
The inverse-diagonal bounds for the positive definite matrix \(H\) bound
\(y,z\) exactly, and the displayed formula recovers \(x\) subject to its
integrality. This is an exhaustive search. A rejected base matrix
includes an explicit square-two vector in its diagnostics.

Testing the full lattice matters. Checking only \(o^\perp\) for
square-two vectors can miss forbidden configurations. Also, one need not
add a separate square-one test: if \(v^2=1\), distinguishedness and
Cauchy--Schwarz imply \((o,v)=\pm1\), and then \((o\mp v)^2=2\).

A square-two vector remains present in every overlattice. It is
therefore safe to reject a base lattice before considering its
saturations.

\section{Enumerating every
saturation}\label{enumerating-every-saturation}

Here is the finite parametrization used in the code. It is particularly
simple because both rank-two labels must remain primitive.

Let \(M=\langle o,u_1,u_2\rangle\) be a retained base lattice, and let
\(N\) be an integral overlattice of finite index \(n\) in which both
labels remain primitive. Since \(\langle o,u_1\rangle\) is primitive in
\(N\), the quotient
\[N/\langle o,u_1\rangle\simeq\mathbb Z.\]

The image of \(M/\langle o,u_1\rangle\) has index \(n\) in this infinite
cyclic group. Choosing its positive generator gives a basis
\((o,u_1,w)\) of \(N\) satisfying
\[w=\frac{a o+b u_1+u_2}{n}.\]

Adding integral multiples of \(o\) and \(u_1\) to \(w\) reduces \(a,b\)
uniquely modulo \(n\). We may therefore assume
\[0\le a,b<n.\]

In this basis, \(u_2=nw-a o-bu_1\). The quotient by the second label has
presentation
\[N/\langle o,u_2\rangle\simeq\mathbb Z^2/\langle(-b,n)\rangle.\]

It is torsion-free precisely when
\[\gcd(b,n)=1.\]

Finally, the determinant-index formula gives
\[\det M=n^2\det N,\]

so \(n^2\) must divide \(\det M\). These statements prove completeness
of a finite enumeration over \(n,a,b\). In particular, the quotient
\(N/M\) is cyclic because primitivity forces it to be cyclic; the
program is not discarding possible noncyclic quotients without
justification.

For each candidate, compute
\[
(o,w)=\frac{3a+\epsilon_1b+\epsilon_2}{n},\qquad
(u_1,w)=\frac{\epsilon_1a+a_1b+t}{n},
\]
\[
w^2=\frac{3a^2+a_1b^2+a_2+2\epsilon_1ab+2\epsilon_2a+2tb}{n^2}.
\]

All three quantities must be integers, and distinguishedness
additionally requires
\[w^2\equiv(o,w)\pmod2.\]

The resulting matrix is again checked for square-two vectors. Positivity
is inherited from the rational quadratic space of \(M\).

\subsection{Checkable inclusion
certificates}\label{checkable-inclusion-certificates}

The inclusion \(M\hookrightarrow N\) is represented by
\[E=\begin{pmatrix}1&0&-a\\0&1&-b\\0&0&n\end{pmatrix}.\]

If \(G_N\) is the returned matrix of \(N\), the certificate satisfies
\[E^TG_NE=G_t,\qquad |\det E|=n.\]

Both label determinants remain \(d_1,d_2\). Primitivity of the first is
visible from the first two basis columns; primitivity of the second is
equivalent to the gcd of its two-by-two minors being one. The tests
verify these identities independently for the emitted certificates.

\subsection{Why isometry deduplication is
necessary}\label{why-isometry-deduplication-is-necessary}

A component parametrizes cubic fourfolds, not a chosen ordered pair of
labels or a chosen basis. The same rank-three lattice may arise from
different values of \(t\), from different saturation descriptions, or
from distinct pairs of rank-two sublattices inside it.

The program consequently identifies matrices only after an exact
integral isometry test preserving \(o\). It seeks a unimodular matrix
\(P\) with
\[P^TG_2P=G_1,\qquad P(1,0,0)^T=(1,0,0)^T.\]

The remaining two columns have prescribed norms and prescribed pairings
with \(o\), so finite coordinate bounds again make this search
exhaustive. Equality of determinants is used only as a quick necessary
check; it is not treated as an isometry criterion. When several origins
produce the same component, their certificates are retained in a common
representative basis.

Different root-free rank-three lattices cannot give one component
properly containing another via a primitive inclusion: a primitive
finite-index inclusion between lattices of the same rank is an
isomorphism. Thus the deduplicated list gives the maximal irreducible
loci needed here.

\section{Worked examples and published
checks}\label{worked-examples-and-published-checks}

The exact regression fixtures include the two requested intersections, a
nontrivial saturation example, and two additional published cases.
Tables of their determinants are collected in the appendix.

\subsection{\texorpdfstring{Five components of
\(\mathcal C_8\cap\mathcal C_{14}\)}{Five components of \textbackslash mathcal C\_8\textbackslash cap\textbackslash mathcal C\_\{14\}}}\label{five-components-of-mathcal-c_8capmathcal-c_14}

The normalized matrix is
\[G_t=\begin{pmatrix}3&1&1\\1&3&t\\1&t&5\end{pmatrix},
\qquad\det G_t=37+2t-3t^2.\]

Positivity allows \(-3\le t\le3\). The end values have square-two
vectors and are excluded. The five retained values are
\(t=-2,-1,0,1,2\), with determinants \(21,32,37,36,29\), respectively.
No additional admissible overlattice contributes a component.

The geometric basis in Auel--Bernardara--Bolognesi--Várilly-Alvarado
{[}ABBVA{]} uses \(T=o+u_2\), so their intersection parameter is
\((P,T)=t+1\). Their parameter range \(-1,0,1,2,3\) therefore agrees
with the normalized range above. A regression test checks the Gram
matrices after this basis conversion, rather than assuming that
parameters used by different authors have identical meanings.

\subsection{\texorpdfstring{Eight components of
\(\mathcal C_8\cap\mathcal C_{26}\)}{Eight components of \textbackslash mathcal C\_8\textbackslash cap\textbackslash mathcal C\_\{26\}}}\label{eight-components-of-mathcal-c_8capmathcal-c_26}

Now
\[G_t=\begin{pmatrix}3&1&1\\1&3&t\\1&t&9\end{pmatrix},
\qquad\det G_t=69+2t-3t^2.\]

The root-free values are \(t=-3,-2,-1,0,1,2,3,4\). Awada {[}A,
Proposition 3.2{]} uses the surface class \(S_{26}=2o+u_2\), giving
parameter \(t+2\) and matrix
\[\begin{pmatrix}3&1&7\\1&3&t+2\\7&t+2&25\end{pmatrix}.\]

Its determinant is unchanged. This explains why a correct implementation
using normalized bases prints different parameter values from the
geometric presentation in the paper.

\subsection{\texorpdfstring{The sixth component of
\(\mathcal C_{20}\cap\mathcal C_{12}\)}{The sixth component of \textbackslash mathcal C\_\{20\}\textbackslash cap\textbackslash mathcal C\_\{12\}}}\label{the-sixth-component-of-mathcal-c_20capmathcal-c_12}

Use the basis order in Yang--Yu {[}YY, Theorem 7.11{]}:
\[M_t=\begin{pmatrix}3&1&0\\1&7&t\\0&t&4\end{pmatrix}.\]

The root-free base lattices have \(0\le t\le4\) and determinants
\(80,77,68,53,32\). At \(t=0\), take
\[w=\frac{o+u_1+u_2}{2}.\]

The basis \((o,u_1,w)\) gives
\[G_N=\begin{pmatrix}3&1&2\\1&7&4\\2&4&4\end{pmatrix},\qquad\det G_N=20.\]

The inclusion has index two. Both original labels remain primitive,
distinguishedness holds, and the new lattice has no square-two vector.
It yields a sixth component. The alternative basis \((o,o-w,u_1)\) gives
\[\begin{pmatrix}3&1&1\\1&3&-3\\1&-3&7\end{pmatrix},\]

exactly the matrix displayed in the published theorem. This example is
the regression that detects omission of saturation.

\section{K3 associations: what is actually being
tested}\label{k3-associations-what-is-actually-being-tested}

For a primitive label \(K_d\), Hassett's untwisted polarized K3
association asks for a polarized K3 surface \((S,L)\), \(L^2=d\), and an
integral Hodge isometry
\[K_d^\perp\simeq L^\perp(-1).\]

The arithmetic condition is
\[4\nmid d,\qquad9\nmid d,\qquad
p\nmid d\text{ for every odd prime }p\equiv2\pmod3.\]

Together with nonemptiness, this is the admissibility criterion for a
label {[}H, Theorem 1.0.2 / 5.1.3{]}. The Tate twist in this expression
changes the sign of the quadratic form.

There are three distinctions that the interface preserves. First, the
component discriminant \(\det A(X)\) is not a label discriminant: one
must look for primitive rank-two sublattices containing \(o\). Second, a
nonadmissible label does not exclude other admissible labels on the same
cubic. For example, a point in \(\mathcal C_8\cap\mathcal C_{14}\) has
the admissible label 14 even though 8 is nonadmissible. Third, a K3
surface equipped with a nonzero Brauer class, often called a twisted K3
surface, is different from the untwisted association above.

\subsection{An exact test for the entire rank-three
lattice}\label{an-exact-test-for-the-entire-rank-three-lattice}

Write the matrix of a component as
\[G=\begin{pmatrix}3&p&q\\p&r&s\\q&s&t\end{pmatrix}.\]

For a vector \(v=x o+y u+z w\), the determinant of
\(\langle o,v\rangle\) is independent of \(x\) and equals
\[f(y,z)=(3r-p^2)y^2+2(3s-pq)yz+(3t-q^2)z^2.\]

The pair spans a primitive rank-two sublattice if and only if
\(\gcd(y,z)=1\). Let \(\lambda\) be the gcd of the three
\emph{polynomial coefficients} of \(f\), including the doubled mixed
coefficient, and write \(f=\lambda g\).

Yang--Yu's Theorems 8.8--8.9 give an exact criterion {[}YY{]}. In the
nonsplit case \(\gcd(3,p,q)=1\), an admissible label exists precisely
when \(\lambda\) passes the prime-divisibility condition above. In the
split case \(\gcd(3,p,q)=3\), one additionally requires that \(g\)
represent \(1\) modulo three. This last condition is checked on the nine
pairs in \((\mathbb Z/3\mathbb Z)^2\).

The content \(\lambda\) is not itself asserted to be a Hassett label. In
particular, \(\lambda=2\) or \(6\) is allowed; imposing \(\lambda>6\)
would be a bug. Nor can the split congruence check be omitted. For
example, \(G=\operatorname{diag}(3,4,6)\) gives \(f=12y^2+18z^2\),
content six, and primitive form \(2y^2+3z^2\), which never represents
one modulo three. It is nonadmissible. In contrast,
\(\operatorname{diag}(3,4,10)\) gives primitive form \(2y^2+5z^2\),
which represents one modulo three at \((1,1)\), and is admissible.

An intersection can reveal an admissible label that neither input
provides. In \(\mathcal C_8\cap\mathcal C_{12}\), the normalized lattice
with \(t=1\) has determinant 29. The primitive vector \(v=u_1+u_2\)
satisfies \(v^2=9\) and \((o,v)=1\), and \(\langle o,v\rangle\) is
primitive because \(\gcd(1,1)=1\). Its discriminant is \(3\cdot9-1=26\).
This component therefore has an untwisted K3 association although both
input labels 8 and 12 are nonadmissible. The components with \(t=0,2\)
instead have content four and are generically nonadmissible.

A positive result persists throughout the component: a primitive
admissible label in \(M\) remains primitive after a primitive inclusion
\(M\hookrightarrow A(X)\). A negative result concerns cubic fourfolds
with \(A(X)=M\), hence the very general point of that component. It must
not be extended to every point: higher-rank specializations can acquire
an admissible label.

The Python and Julia function \texttt{k3\_certificate} and Macaulay2's
\texttt{k3Certificate} make this decision constructive. An affirmative
answer contains a primitive vector and its rank-two Gram matrix; a
negative answer contains the arithmetic obstruction. Increasing
discriminant bounds exhaust the witness search, under one shared
resource budget. The older bounded witness functions can return no
result within their bound even when an admissible marking exists.
Resource exhaustion is never presented as a negative mathematical
result.

\section{Rationality: geometric constructions and the limits of the
output}\label{rationality-geometric-constructions-and-the-limits-of-the-output}

The lattice calculation identifies algebraic Hodge conditions. It does
not produce a birational map to projective space. The historical
Hassett/Kuznetsov picture links untwisted K3 associations, geometric K3
categories, and rationality; this project does not turn that picture
into an unconditional numerical equivalence.

There are concrete positive results. Every smooth cubic in
\(\mathcal C_{14}\) is rational by Bolognesi--Russo--Staglianò {[}BRS,
Theorems 2.9 and 3.4{]}, and also by {[}RS, Theorem 2{]}. Every smooth
cubic in \(\mathcal C_{26}\) is rational by {[}RS, Theorem 4{]}. These
geometric results apply to every member of both requested plane
intersections. They are stronger than a statement about a general member
of a divisor, which alone would not establish rationality on every
intersection component.

For a cubic containing a plane \(P\), projection from the plane gives a
quadric surface fibration over \(\mathbb P^2\). The residual quadric
class is \(Q=o-P\). An algebraic class with odd intersection with \(Q\)
supplies the classical rational-section criterion, hence rationality;
with smooth discriminant, the same condition is expressed by vanishing
of the Clifford Brauer class on the degree-two K3 discriminant cover
{[}A, Section 3.1{]}.

In the normalized \(\mathcal C_8\cap\mathcal C_{14}\) matrix,
\[Q\cdot o=2,\qquad Q\cdot P=-2,\qquad Q\cdot u_2=1-t.\]

Thus an odd pairing occurs exactly when \(t\) is even, assuming this
rank-three lattice is the full algebraic lattice. More explicitly,
\((t/2)o+u_2\) pairs to one with \(Q\) for even \(t\). For \(t=-1\), the
determinant is 32 and this section criterion fails. Under the good-plane
hypothesis, the Clifford invariant is nontrivial {[}ABBVA, Theorem A{]},
although the cubic is rational by {[}BRS{]}. We use this later
rationality proof: BRS Section 3 addresses a gap involving openness in
the earlier pfaffian argument, and Theorem 3.7 shows that the pfaffian
locus is not open in \(\mathcal C_{14}\). Failure of one rationality
construction is therefore not an irrationality proof.

For \(\mathcal C_8\cap\mathcal C_{26}\), Awada's parameter is \(t+2\),
so her even-parameter criterion agrees with even \(t\) {[}A, Theorem
3.3{]}. The statement about the Brauer class requires a good plane and
the specified full rank-three algebraic lattice. An existing section
class persists at higher-rank specializations; failure of the criterion
at rank three need not persist.

Two components of \(\mathcal C_{12}\cap\mathcal C_{20}\) also have
unconditional rationality conclusions. In the normalized order
\((20,12)\), the primitive vector \(u_1-u_2\) has label \(26\) at
\(t=1\) and label \(14\) at \(t=3\). The entire corresponding loci
therefore lie in \(\mathcal C_{26}\) and \(\mathcal C_{14}\),
respectively. The other four component lattices have a content-four
obstruction to untwisted K3 admissibility. This computation does not
decide their rationality.

\subsection{The categorical bridge}\label{the-categorical-bridge}

The derived category has a semiorthogonal decomposition
\[D^b(X)=\langle\operatorname{Ku}(X),\mathcal O_X,\mathcal O_X(1),\mathcal O_X(2)\rangle.\]

The residual category \(\operatorname{Ku}(X)\) is the Kuznetsov
component. Its Serre functor is the shift by two, so it is a
two-Calabi--Yau category, with the categorical structure characteristic
of a K3 surface. Kuznetsov's 2010 Conjecture 1.1 proposes that \(X\) is
rational precisely when \(\operatorname{Ku}(X)\) is equivalent to
\(D^b(S)\) for an untwisted K3 surface \(S\) {[}K{]}. An equivalence
with a twisted category \(D^b(S,\alpha)\), where \(\alpha\ne0\), is a
different condition and does not by itself meet that conjectural
criterion.

Addington--Thomas established generic agreement between Hassett's
Hodge-theoretic K3 condition and the existence of an untwisted geometric
K3 category {[}AT{]}. This historical result explains why the arithmetic
admissibility test belongs in a discussion of rationality. The explorer
computes the lattice condition; it does not construct a derived
equivalence, compute the Kuznetsov component, or establish the
conjectured equivalence with rationality.

\subsection{Current context, checked 14 November
2026}\label{current-context-checked-14-november-2026}

The older statement that no irrational smooth complex cubic fourfold is
known is no longer appropriate. Katzarkov--Kontsevich--Pantev--Yu prove
that a very general cubic fourfold is irrational using Hodge atoms
{[}KKPY, current arXiv version dated 6 March 2026{]}. Guéré's September
2026 version proves that rationality implies an isomorphism
\[H^4(X,\mathbb Q)_{\mathrm{prim}}\simeq H^2(S,\mathbb Q)(-1)\]

for a projective K3 surface \(S\) {[}G, Theorem 63{]}. This statement is
explicitly about rational Hodge structures. It should not silently be
substituted for an integral polarized Hassett association: the integral
lattice and its pairing contain extra information. The twist here is a
Tate twist, not a Brauer twist. These arXiv preprints provide context;
the program invokes neither result in its correctness argument and
derives no additional rationality or irrationality decision from them.

\section{Using and interpreting the
explorer}\label{using-and-interpreting-the-explorer}

The offline explorer displays a finite catalogue generated by the same
Python engine. Selecting two divisors exposes the rank-three components,
and selecting a component shows its matrix, determinant, origins, and K3
test. The candidate scan explains which pairings survive the root test
and where saturation adds a component. JSON exported by the command-line
tool can be loaded into the explorer for additional discriminants.

The network is an incidence schematic. Distances, angles, edge
crossings, and vertex positions have no meaning in the 20-dimensional
moduli space. An edge records the computed component count for a bundled
pair. An unbundled pair is not asserted to have empty intersection. The
display is not a projection of the period domain, and it does not
display intersection multiplicities.

For computational reproducibility, retain the JSON alongside the version
of the code. The matrices and inclusion certificates permit elementary
independent checks; the completeness argument is the finite enumeration
proved above plus the cited geometric theorem. Passing the published
fixtures is a regression check, not a substitute for that argument.

\clearpage
\section{Appendix: validation
targets}\label{appendix-validation-targets}

The table records the published component counts and the sorted
discriminants of their generic rank-three algebraic lattices. The
reproducible validation report in the distribution records the actual
run of the implementation.

\begin{longtable}[]{@{}
  >{\raggedright\arraybackslash}p{(\linewidth - 6\tabcolsep) * \real{0.18}}
  >{\raggedleft\arraybackslash}p{(\linewidth - 6\tabcolsep) * \real{0.17}}
  >{\raggedright\arraybackslash}p{(\linewidth - 6\tabcolsep) * \real{0.43}}
  >{\raggedright\arraybackslash}p{(\linewidth - 6\tabcolsep) * \real{0.22}}@{}}
\toprule\noalign{}
\begin{minipage}[b]{\linewidth}\raggedright
Intersection
\end{minipage} & \begin{minipage}[b]{\linewidth}\raggedleft
Components
\end{minipage} & \begin{minipage}[b]{\linewidth}\raggedright
Sorted component discriminants
\end{minipage} & \begin{minipage}[b]{\linewidth}\raggedright
Primary reference
\end{minipage} \\
\midrule\noalign{}
\endhead
\bottomrule\noalign{}
\endlastfoot
\(\mathcal C_8\cap\mathcal C_{14}\) & 5 & 21, 29, 32, 36, 37 &
{[}ABBVA{]}, Theorem A / 4 \\
\(\mathcal C_8\cap\mathcal C_{26}\) & 8 & 29, 36, 48, 53, 61, 64, 68, 69
& {[}A{]}, Proposition 3.2 \\
\(\mathcal C_8\cap\mathcal C_{38}\) & 10 & 36, 45, 61, 68, 80, 85, 93,
96, 100, 101 & {[}A{]}, Theorem 1.1(2) \\
\(\mathcal C_{20}\cap\mathcal C_{12}\) & 6 & 20, 32, 53, 68, 77, 80 &
{[}YY{]}, Theorem 7.11 \\
\(\mathcal C_{20}\cap\mathcal C_{38}\) & 17 & 45, 77, 92, 120, 133, 157,
168, 188, 197, 213, 220, 232, 237, 245, 248, 252, 253 & {[}YY{]},
Theorem 7.9 \\
\end{longtable}

The automated checks also cover input rejection, swapping the two
discriminants, exact isometry, explicit root witnesses, primitive
inclusion certificates, the index-two example above, and the
split/nonsplit K3 criterion. The source release adds the general Hermite
audit and independent verification of constructive K3 certificates.
Current test counts and source fingerprints are recorded by the release
validator rather than fixed in this guide. Live-browser rendering,
downloads, and clipboard behavior remain separate checks.

\clearpage
\section{References}\label{references}

\textbf{{[}BRS{]}} Michele Bolognesi, Francesco Russo, and Giovanni
Staglianò, \emph{Some loci of rational cubic fourfolds}, arXiv version
6, 1 June 2018. Theorems 2.9, 3.4 and 3.7.
\href{https://arxiv.org/abs/1504.05863v6}{arXiv:1504.05863v6}.

\textbf{{[}H{]}} Brendan Hassett, \emph{Special cubic fourfolds},
Compositio Mathematica \textbf{120} (2000), 1--23. Theorems 1.0.1, 1.0.2
and 5.1.3.
\href{https://doi.org/10.1023/A:1001706324425}{doi:10.1023/A:1001706324425}.

\textbf{{[}YY{]}} Song Yang and Xun Yu, \emph{On lattice polarizable
cubic fourfolds}, Research in the Mathematical Sciences \textbf{10}
(2023), article 2. Algorithm 7.6; Theorems 1.1, 7.9, 7.11, 8.8 and 8.9.
\href{https://arxiv.org/abs/2103.09132}{arXiv:2103.09132},
\href{https://doi.org/10.1007/s40687-022-00368-6}{doi:10.1007/s40687-022-00368-6}.
Numbering follows the available arXiv version.

\textbf{{[}ABBVA{]}} Asher Auel, Marcello Bernardara, Michele Bolognesi
and Anthony Várilly-Alvarado, \emph{Cubic fourfolds containing a plane
and a quintic del Pezzo surface}, Algebraic Geometry \textbf{1} (2014),
181--193. \href{https://arxiv.org/abs/1205.0237v4}{arXiv:1205.0237v4}.

\textbf{{[}A{]}} Hanine Awada, \emph{Rational fibered cubic fourfolds},
arXiv version 2, 9 July 2024. Theorem 1.1 and Proposition 3.2.
\href{https://arxiv.org/abs/1910.10182v2}{arXiv:1910.10182v2}. Earlier
versions used the title \emph{Rational fibered cubic fourfolds with
nontrivial Brauer classes}.

\textbf{{[}RS{]}} Francesco Russo and Giovanni Staglianò,
\emph{Congruences of 5-secant conics and the rationality of some
admissible cubic fourfolds}, Duke Mathematical Journal \textbf{168}
(2019), 849--865.
\href{https://arxiv.org/abs/1707.00999}{arXiv:1707.00999},
\href{https://doi.org/10.1215/00127094-2018-0053}{doi:10.1215/00127094-2018-0053}.

\textbf{{[}K{]}} Alexander Kuznetsov, \emph{Derived categories of cubic
fourfolds}, in \emph{Cohomological and Geometric Approaches to
Rationality Problems}, Progress in Mathematics \textbf{282} (2010),
219--243. Conjecture 1.1.
\href{https://arxiv.org/abs/0808.3351}{arXiv:0808.3351}.

\textbf{{[}AT{]}} Nicolas Addington and Richard Thomas, \emph{Hodge
theory and derived categories of cubic fourfolds}, Duke Mathematical
Journal \textbf{163} (2014), 1885--1927.
\href{https://arxiv.org/abs/1211.3758}{arXiv:1211.3758}.

\textbf{{[}KKPY{]}} Ludmil Katzarkov, Maxim Kontsevich, Tony Pantev and
Tony Yue Yu, \emph{Birational invariants from Hodge structures and
quantum multiplication}, version 2, 6 March 2026.
\href{https://arxiv.org/abs/2508.05105v2}{arXiv:2508.05105v2}.

\textbf{{[}G{]}} Jérémy Guéré, \emph{On the irrationality of cubic
fourfolds}, version 2, 4 September 2026. Theorem 63.
\href{https://arxiv.org/abs/2603.04518v2}{arXiv:2603.04518v2}.

\textbf{Authorship and provenance.} This software and explanatory text
were developed with AI assistance. The mathematical algorithm and
geometric theorems are credited to the authors above. The
implementation, saturation parametrization proof, certificates, and
tests should be reviewed as research software; the project is not a
formal proof assistant verification.

\end{document}
